% !TEX root = ../main.tex
\section{Motivación y beneficios esperados}
\label{sec:motivacion}
El escriba clínico humano registra información del encuentro asistencial para facilitar su incorporación a la historia clínica electrónica. Los escribas digitales buscan automatizar parte de esa tarea mediante la captura de la conversación y la generación de un borrador. Su propuesta responde a una carga documental relevante: un estudio de observación de 57 médicos de cuatro especialidades en Estados Unidos atribuyó el 49,2\% de la jornada en consulta al trabajo con la historia electrónica y de escritorio~\cite{7MXFHMBA,sinsky2016}.

\textbf{Adopción y mercado.} PHTI describió en 2025 una expansión excepcionalmente rápida, basada en la experiencia de directivos y empresas del sector. Esta valoración respalda el interés por la tecnología, pero no establece un récord histórico mundial de adopción. La revisión de mercado del presente informe reúne 43 proveedores con diferentes modelos comerciales y grados de integración; su amplitud justifica una selección orientada a los requisitos locales~\cite{FKHM4UFL}.

\textbf{Beneficio esperado.} Una revisión de 35 estudios identifica resultados generalmente favorables o mixtos sobre bienestar, con predominio de investigaciones estadounidenses y medidas percibidas~\cite{4PZBAAZ7}. En el ensayo aleatorizado de UCLA, Nabla redujo el tiempo de documentación un 9,5\% frente al control; DAX no mostró una diferencia significativa~\cite{PR93N8FN}. Por tanto, el ahorro debe estimarse localmente, sin adoptar un porcentaje único como expectativa de todos los productos.

\textbf{Incertidumbres.} Persisten preguntas sobre errores, variación entre especialidades, efectos sobre el razonamiento profesional y cambios acumulados en la historia clínica. La posible dependencia cognitiva es una hipótesis de investigación, no un perjuicio demostrado de estos sistemas~\cite{ZI87XIWC}. El objetivo del proyecto es comprobar si el alivio documental se traduce en una mejora útil y sostenible para el hospital.
